% GENERATED FILE. Edit canonical lessons, then run npm run generate:books.
% canonical-source-hash: fnv1a64:dc1f108d5d5f771e
% canonical-lessons: GE-C309-schmuecken, GE-C309-girlande, GE-C309-aufhaengen, GE-C309-luftballon, GE-C309-feuerwerk, GE-R309-review-302-309

\chapter{Alles ist geschmückt}
\label{ch:ge-alles-ist-geschmuckt}
\hlchaptermodality{309}

\begin{chapteropening}
\textbf{By the end of this chapter:} \emph{I can describe decorations: to decorate, a garland, to hang up, a balloon and fireworks.}

Complete the last lesson of chapter 309: I can describe decorations: to decorate, a garland, to hang up, a balloon and fireworks.
\end{chapteropening}

\section[schmücken]{schmücken — to decorate, to adorn}

\label{lesson:GE-C309-schmuecken}



\begin{quote}
Before the new one: say the German for the applause, then the German for to applaud.
\end{quote}

\subsection*{You'll want to know: schmücken}
\textbf{schmücken} — \textquotedblleft{}to decorate, to adorn\textquotedblright{}.

\begin{quote}
Wir schmückten den Saal mit Blumen.
\end{quote}

\begin{quote}
\emph{We decorated the hall with flowers.}
\end{quote}

\emph{Der Schmuck} — jewellery or decoration.

\begin{cousinweb}
\emph{schmücken} is usually traced to Middle High German \emph{smücken}, which EtymDB glosses \textquotedblleft{}to dress, to adorn\textquotedblright{}.
\end{cousinweb}

\begin{grammarlens}[title={What you've built}]
One more word for telling, explaining and describing at B1.
\end{grammarlens}

\subsection*{Your turn}
\begin{itemize}
\raggedright
  \item \emph{Say it:} \emph{schmücken}
  \item \emph{Say it:} \emph{schmücken}, once more
  \item \emph{Say it:} a sentence of your own with \emph{schmücken}
  \item \emph{Recall:} say the German for the applause, then the German for to applaud, then say \emph{schmücken} again
\end{itemize}

\begin{tcolorbox}[breakable,colback=teal!4,colframe=teal!35!black,title={Before you move on}]
What does schmücken mean? (\textquotedblleft{}To decorate, to adorn\textquotedblright{}.) Say it once more.
\end{tcolorbox}
\section[die Girlande]{die Girlande — the garland, the paper chain}

\label{lesson:GE-C309-girlande}



\begin{quote}
Before the new one: say the German for to applaud, then the German for to decorate.
\end{quote}

\subsection*{You'll want to know: die Girlande}
\textbf{die Girlande} — \textquotedblleft{}the garland, the paper chain\textquotedblright{}.

\begin{quote}
Über der Tür hing eine bunte Girlande.
\end{quote}

\begin{quote}
\emph{A colourful garland hung over the door.}
\end{quote}

Plural: \emph{die Girlanden}.

\begin{cousinweb}
\emph{Girlande} was borrowed from French \emph{guirlande} \textquotedblleft{}garland\textquotedblright{}, which came from Italian \emph{ghirlanda} (EtymDB).
\end{cousinweb}

\begin{grammarlens}[title={What you've built}]
One more word for telling, explaining and describing at B1.
\end{grammarlens}

\subsection*{Your turn}
\begin{itemize}
\raggedright
  \item \emph{Say it:} \emph{die Girlande}
  \item \emph{Say it:} \emph{die Girlande}, once more
  \item \emph{Say it:} a sentence of your own with \emph{die Girlande}
  \item \emph{Recall:} say the German for to applaud, then the German for to decorate, then say \emph{die Girlande} again
\end{itemize}

\begin{tcolorbox}[breakable,colback=teal!4,colframe=teal!35!black,title={Before you move on}]
What does die Girlande mean? (\textquotedblleft{}The garland, the paper chain\textquotedblright{}.) Say it once more.
\end{tcolorbox}
\section[aufhängen]{aufhängen — to hang up}

\label{lesson:GE-C309-aufhaengen}



\begin{quote}
Before the new one: say the German for to decorate, then the German for the garland.
\end{quote}

\subsection*{You'll want to know: aufhängen}
\textbf{aufhängen} — \textquotedblleft{}to hang up\textquotedblright{}.

\begin{quote}
Die Kinder hängten Bilder an der Wand auf.
\end{quote}

\begin{quote}
\emph{The children hung pictures on the wall.}
\end{quote}

Separable: \emph{er hängt auf}. On the phone, to hang up.

\begin{cousinweb}
\emph{aufhängen} joins \emph{auf} \textquotedblleft{}up\textquotedblright{} and \emph{hängen} \textquotedblleft{}to hang\textquotedblright{}, an inherited Germanic verb and a cousin of English \emph{hang}; EtymDB lists no parent for the compound.
\end{cousinweb}

\begin{grammarlens}[title={What you've built}]
One more word for telling, explaining and describing at B1.
\end{grammarlens}

\subsection*{Your turn}
\begin{itemize}
\raggedright
  \item \emph{Say it:} \emph{aufhängen}
  \item \emph{Say it:} \emph{aufhängen}, once more
  \item \emph{Say it:} a sentence of your own with \emph{aufhängen}
  \item \emph{Recall:} say the German for to decorate, then the German for the garland, then say \emph{aufhängen} again
\end{itemize}

\begin{tcolorbox}[breakable,colback=teal!4,colframe=teal!35!black,title={Before you move on}]
What does aufhängen mean? (\textquotedblleft{}To hang up\textquotedblright{}.) Say it once more.
\end{tcolorbox}
\section[der Luftballon]{der Luftballon — the balloon}

\label{lesson:GE-C309-luftballon}



\begin{quote}
Before the new one: say the German for the garland, then the German for to hang up.
\end{quote}

\subsection*{You'll want to know: der Luftballon}
\textbf{der Luftballon} — \textquotedblleft{}the balloon\textquotedblright{}.

\begin{quote}
Jedes Kind bekam einen roten Luftballon.
\end{quote}

\begin{quote}
\emph{Every child got a red balloon.}
\end{quote}

Plural: \emph{die Luftballons}.

\begin{cousinweb}
\emph{Luftballon} joins \emph{Luft} \textquotedblleft{}air\textquotedblright{}, Old High German \emph{luft} (EtymDB), and \emph{Ballon}, borrowed from French \emph{ballon} \textquotedblleft{}a large ball\textquotedblright{}.
\end{cousinweb}

\begin{grammarlens}[title={What you've built}]
One more word for telling, explaining and describing at B1.
\end{grammarlens}

\subsection*{Your turn}
\begin{itemize}
\raggedright
  \item \emph{Say it:} \emph{der Luftballon}
  \item \emph{Say it:} \emph{der Luftballon}, once more
  \item \emph{Say it:} a sentence of your own with \emph{der Luftballon}
  \item \emph{Recall:} say the German for the garland, then the German for to hang up, then say \emph{der Luftballon} again
\end{itemize}

\begin{tcolorbox}[breakable,colback=teal!4,colframe=teal!35!black,title={Before you move on}]
What does der Luftballon mean? (\textquotedblleft{}The balloon\textquotedblright{}.) Say it once more.
\end{tcolorbox}
\section[das Feuerwerk]{das Feuerwerk — the fireworks, the firework display}

\label{lesson:GE-C309-feuerwerk}



\begin{quote}
Before the new one: say the German for to hang up, then the German for the balloon.
\end{quote}

\subsection*{You'll want to know: das Feuerwerk}
\textbf{das Feuerwerk} — \textquotedblleft{}the fireworks, the firework display\textquotedblright{}.

\begin{quote}
Um Mitternacht gab es ein großes Feuerwerk.
\end{quote}

\begin{quote}
\emph{At midnight there was a big firework display.}
\end{quote}

Plural: \emph{die Feuerwerke}. One rocket is \emph{eine Rakete}.

\begin{cousinweb}
\emph{Feuerwerk} joins \emph{Feuer} \textquotedblleft{}fire\textquotedblright{}, a cousin of English \emph{fire} (EtymDB), and \emph{Werk} \textquotedblleft{}work\textquotedblright{}, both inherited Germanic words.
\end{cousinweb}

\begin{grammarlens}[title={What you've built}]
One more word for telling, explaining and describing at B1.
\end{grammarlens}

\subsection*{Your turn}
\begin{itemize}
\raggedright
  \item \emph{Say it:} \emph{das Feuerwerk}
  \item \emph{Say it:} \emph{das Feuerwerk}, once more
  \item \emph{Say it:} a sentence of your own with \emph{das Feuerwerk}
  \item \emph{Recall:} say the German for to hang up, then the German for the balloon, then say \emph{das Feuerwerk} again
\end{itemize}

\begin{tcolorbox}[breakable,colback=teal!4,colframe=teal!35!black,title={Before you move on}]
What does das Feuerwerk mean? (\textquotedblleft{}The fireworks, the firework display\textquotedblright{}.) Say it once more.
\end{tcolorbox}
\section[(dialogue)]{Words from chapters 302–309 — a review}

\label{lesson:GE-R309-review-302-309}



\begin{quote}
Say the German for the pension and the retirement.
\end{quote}

\subsection*{Your turn}
Say these aloud:

\begin{itemize}
\raggedright
  \item the pension, the retirement, the pensioner, the anniversary, the farewell
  \item the death, to bury, the funeral, the coffin, the cemetery
  \item the grief, to mourn, the condolences, the widow, to inherit
  \item the celebration, solemn, the public holiday, the ceremony, to open officially
  \item to congratulate, the good wishes, to present, to give presents to, to unwrap
  \item the host, to treat guests, the guest of honour, the speech, the toast
  \item to clink glasses, the sparkling wine, to cheer, the applause, to applaud
  \item to decorate, the garland, to hang up, the balloon, the fireworks
\end{itemize}

\begin{tcolorbox}[breakable,colback=teal!4,colframe=teal!35!black,title={Before you move on}]
The pension? (\textbf{die Rente}.) The retirement? (\textbf{der Ruhestand}.) The pensioner? (\textbf{der Rentner}.) The anniversary? (\textbf{das Jubiläum}.) The farewell? (\textbf{der Abschied}.) The death? (\textbf{der Tod}.) To bury? (\textbf{beerdigen}.) The funeral? (\textbf{die Beerdigung}.) The coffin? (\textbf{der Sarg}.) The cemetery? (\textbf{der Friedhof}.) The grief? (\textbf{die Trauer}.) To mourn? (\textbf{trauern}.) The condolences? (\textbf{das Beileid}.) The widow? (\textbf{die Witwe}.) To inherit? (\textbf{erben}.) The celebration? (\textbf{die Feier}.) Solemn? (\textbf{feierlich}.) The public holiday? (\textbf{der Feiertag}.) The ceremony? (\textbf{die Zeremonie}.) To open officially? (\textbf{einweihen}.) To congratulate? (\textbf{gratulieren}.) The good wishes? (\textbf{der Glückwunsch}.) To present? (\textbf{überreichen}.) To give presents to? (\textbf{beschenken}.) To unwrap? (\textbf{auspacken}.) The host? (\textbf{der Gastgeber}.) To treat guests? (\textbf{bewirten}.) The guest of honour? (\textbf{der Ehrengast}.) The speech? (\textbf{die Ansprache}.) The toast? (\textbf{der Trinkspruch}.) To clink glasses? (\textbf{anstoßen}.) The sparkling wine? (\textbf{der Sekt}.) To cheer? (\textbf{jubeln}.) The applause? (\textbf{der Beifall}.) To applaud? (\textbf{applaudieren}.) To decorate? (\textbf{schmücken}.) The garland? (\textbf{die Girlande}.) To hang up? (\textbf{aufhängen}.) The balloon? (\textbf{der Luftballon}.) The fireworks? (\textbf{das Feuerwerk}.)
\end{tcolorbox}
